% Part: first-order-logic
% Chapter: natural-deduction
% Section: proof-theoretic-notions

\documentclass[../../../include/open-logic-section]{subfiles}

\begin{document}

\iftag{FOL}
      {\olfileid{fol}{ntd}{ptn}}
      {\olfileid{pl}{ntd}{ptn}}

\olsection{सिद्धता-उपपत्तीय संकल्पना}

\begin{editorial}
या विभागात नैसर्गिक निगमनासाठी सिद्धतायोग्यता संबंध आणि सुसंगतता यांच्या
व्याख्या एकत्रित केल्या आहेत.
\end{editorial}

\begin{explain}
जशा आपण अनेक महत्त्वाच्या चिन्हार्थविषयक संकल्पना (वैधता, तार्किक निष्पन्नता,
पूर्ततायोग्यता) परिभाषित केल्या आहेत, तशाच आता त्यांना अनुरूप
\emph{सिद्धता-उपपत्तीय संकल्पना} परिभाषित करू. या संकल्पना !!{structure}s मध्ये
!!{sentence}s ची पूर्ति होते की नाही याच्या आधारे परिभाषित केलेल्या नसून,
काही !!{sentence}s ची इतरांपासूनची !!{derivability} किंवा
!!{nonderivability} यांचा आधार घेऊन परिभाषित केल्या आहेत. या दोन प्रकारच्या
संकल्पना एकमेकांशी जुळतात, हा एक महत्त्वाचा शोध होता. त्या जुळतात, हाच
\emph{निर्दोषता प्रमेय} आणि \emph{संपूर्णता प्रमेय} यांचा आशय आहे.
\end{explain}


\begin{defn}[प्रमेये]
एखादे !!{sentence}~$!A$ हे \emph{प्रमेय} असते, जर नैसर्गिक निगमनात $!A$ ची
अशी !!a{derivation} असेल की तिच्यातील सर्व गृहीतके !!{discharged} आहेत.
$!A$ हे प्रमेय असेल तर आपण $\Proves !A$ लिहितो आणि ते प्रमेय नसेल तर
$\Proves/ !A$ लिहितो.
\end{defn}

\begin{defn}[!!^{derivability}]
!!^a{sentence} $!A$ हे !!{sentence}s च्या संच~$\Gamma$ \emph{पासून
!!{derivable}} असते, म्हणजे $\Gamma \Proves !A$, तेव्हा आणि केवळ तेव्हाच,
जेव्हा निष्कर्ष~$!A$ असणारी अशी !!{derivation} असेल की तिच्यातील प्रत्येक
गृहीतक एकतर !!{discharged} आहे किंवा $\Gamma$ मध्ये आहे. $!A$ हे $\Gamma$
पासून !!{derivable} नसेल, तर आपण $\Gamma \Proves/ !A$ लिहितो.
\end{defn}

\begin{defn}[सुसंगतता]
!!{sentence}s चा संच~$\Gamma$ हा \emph{विसंगत} असतो तेव्हा आणि केवळ तेव्हाच,
जेव्हा $\Gamma \Proves \lfalse$. $\Gamma$ विसंगत नसेल, म्हणजे
$\Gamma \Proves/ \lfalse$ असेल, तर त्याला \emph{सुसंगत} म्हणतो.
\end{defn}

\begin{prop}[परावर्तिता]
\ollabel{prop:reflexivity}
$!A \in \Gamma$ असेल, तर $\Gamma \Proves !A$.
\end{prop}

\begin{proof}
केवळ $!A$ हे गृहीतक स्वतःच $!A$ ची अशी !!a{derivation} आहे की तिच्यातील
प्रत्येक !!{undischarged} गृहीतक (म्हणजे $!A$) $\Gamma$ मध्ये आहे.
\end{proof}


\begin{prop}[एकस्वनिकता]
\ollabel{prop:monotonicity}
$\Gamma \subseteq \Delta$ आणि $\Gamma \Proves !A$ असेल, तर $\Delta
\Proves !A$.
\end{prop}

\begin{proof}
$\Gamma$ पासून $!A$ ची कोणतीही !!{derivation} ही $\Delta$ पासून $!A$ चीही
!!a{derivation} असते.
\end{proof}

\begin{prop}[संक्रमकता]
\ollabel{prop:transitivity}
$\Gamma \Proves !A$ आणि $\{!A\} \cup \Delta \Proves
!B$ असेल, तर $\Gamma \cup \Delta \Proves !B$.
\end{prop}

\begin{proof}
$\Gamma \Proves !A$ असेल, तर $!A$ ची अशी !!a{derivation}~$\delta_0$ असते
की तिच्यातील सर्व !!{undischarged} गृहीतके $\Gamma$ मध्ये असतात. $\{!A\}
\cup \Delta \Proves !B$ असेल, तर $!B$ ची अशी !!a{derivation}~$\delta_1$
असते की तिच्यातील सर्व !!{undischarged} गृहीतके $\{!A\} \cup \Delta$ मध्ये
असतात. आता पुढील रचना विचारात घ्या:
\begin{prooftree}
  \AxiomC{$\Delta, \Discharge{!A}{1}$}
  \RightLabel{$\delta_1$}
  \DeduceC{$!B$}
  \DischargeRule{\Intro{\lif}}{1}
  \UnaryInfC{$!A \lif !B$}
  \AxiomC{$\Gamma$}
  \RightLabel{$\delta_0$}
  \DeduceC{$!A$}
  \RightLabel{\Elim{\lif}}
  \BinaryInfC{$!B$}
\end{prooftree}
आता सर्व !!{undischarged} गृहीतके $\Gamma \cup \Delta$ मध्ये आहेत; म्हणून
यावरून $\Gamma \cup \Delta \Proves !B$ हे दिसते.
\end{proof}

$\Gamma = \{!A_1, !A_2, \ldots, !A_k\}$ हा सांत संच असेल, तेव्हा
$\Gamma \Proves !B$ ऐवजी आपण $!A_1,!A_2,\ldots,!A_k \Proves !B$ हे
सरलीकृत संकेतलेखन वापरू शकतो. विशेषतः, $!A \Proves !B$ याचा अर्थ
$\{!A\} \Proves !B$ असा होतो.

लक्षात घ्या की $\Gamma \Proves !A$ आणि $!A
\Proves !B$ असेल, तर $\Gamma \Proves !B$. तसेच $!A_1,
\dots, !A_n \Proves !B$ असेल आणि प्रत्येक~$i$ साठी $\Gamma \Proves !A_i$
असेल, तर $\Gamma \Proves !B$, हेसुद्धा यावरून मिळते.

\begin{prop}
\ollabel{prop:incons}
पुढील विधाने सममूल्य आहेत.
    \begin{enumerate}
      \item \( \Gamma \) विसंगत आहे.
      \item प्रत्येक !!{sentence}~\( {!A} \) साठी \( \Gamma \Proves {!A} \).
      \item एखाद्या !!{sentence}~\( {!A} \) साठी \( \Gamma \Proves {!A} \) आणि \( \Gamma \Proves \lnot {!A} \).
    \end{enumerate}
\end{prop}

\begin{proof}
सराव.
\end{proof}

\tagprob{FOL}
\begin{prob}
\olref[fol][ntd][ptn]{prop:incons} सिद्ध करा.
\end{prob}
\tagendprob

\tagprob{notFOL}
\begin{prob}
\olref[pl][ntd][ptn]{prop:incons} सिद्ध करा.
\end{prob}
\tagendprob

\begin{prop}[संहतता]
  \ollabel{prop:proves-compact}
  \begin{enumerate}
  \item $\Gamma \Proves !A$ असेल, तर एखादा सांत उपसंच $\Gamma_0
    \subseteq \Gamma$ असा आहे की $\Gamma_0 \Proves !A$.
  \item $\Gamma$ चा प्रत्येक सांत उपसंच सुसंगत असेल, तर $\Gamma$ सुसंगत
    असतो.
  \end{enumerate}
\end{prop}

\begin{proof}
  \begin{enumerate}
    \item $\Gamma \Proves !A$ असेल, तर $\Gamma$ पासून $!A$ ची
      !!a{derivation}~$\delta$ असते. $\Gamma_0$ हा $\delta$ मधील
      !!{undischarged} गृहीतकांचा संच असू द्या. प्रत्येक !!{derivation} सांत
      असल्यामुळे $\Gamma_0$ मधील !!{sentence}s ची संख्या सांतच असते.
      म्हणून $\delta$ ही सांत~$\Gamma_0 \subseteq \Gamma$ पासून $!A$ ची
      !!a{derivation} आहे.
    \item (1) मध्ये $!A
      \ident \lfalse$ हे विशेष प्रकरण घेतल्यावर मिळणारे हे निषेधात्मक उलट
      विधान आहे.
  \end{enumerate}
\end{proof}

\end{document}
